\section{Structural-Recovery Methods (Study~1)}
\label{app:methods}

This appendix specifies the structural-recovery study (Study~1)
and the source identifiers used in its preserved records.
The fitted operators act in training-PCA coordinates without
second-moment whitening.
The coordinate study (Study~2) uses separate banks and
training-energy-normalized prediction losses
(\cref{app:characterization,app:pending-protocols}).
Its prediction scores are not directly comparable with the
per-target relative errors defined here.
The self-whitened class risks in \cref{app:spectral-risk} and
the centering analysis in
\cref{app:template-centering-sensitivity} have their own
observation and evaluation contracts.

\paragraph{Names in archived records.}
The descriptive names in the main text map to the unchanged
source identifiers below.
These identifiers locate the corresponding code and results;
they do not denote additional independent datasets.
\begin{center}
\small
\begin{tabular}{@{}p{.36\textwidth}p{.12\textwidth}p{.44\textwidth}@{}}
\toprule
Main-text comparison & Source name & Scope \\
\midrule
All-element structural recovery
& FULL
& Independent and joint fits, prediction and consistency
  diagnostics, and controls. \\
Generator-only comparison
& AFS-A
& Generator fits evaluated on excluded action elements
  and word controls. \\
Group-label registration control
& AFS-B
& Fixed fitted operators evaluated under non-automorphic
  group-element relabeling. \\
Training-correspondence controls
& AFS-C
& Ridge refits under template-wise state relabeling
  or target scrambling. \\
Rank-matched lexical comparison
& AFS-D
& Lexical and internal readouts matched in training projected
  rank; input spectra and the PCA procedure support the match. \\
\bottomrule
\end{tabular}
\end{center}

\subsection{States, actions, and prompts}

The $C_7$ states are ordered Monday, Tuesday, Wednesday,
Thursday, Friday, Saturday, Sunday, indexed $0,\ldots,6$.
The generator $r$ applies $i\mapsto i+1\pmod7$.
This action is free and transitive.
All seven group elements are available for FULL
element-specific fitting; the focal AFS-A fit uses only
the one-step generator.

The $D_4$ states are NW, NE, SE, SW, indexed $0,1,2,3$.
Rotation is clockwise, $r(i)=i+1\pmod4$.
The generator $s$ is reflection across the vertical axis:
\[
 s(0)=1,\qquad s(1)=0,\qquad s(2)=3,\qquad s(3)=2.
\]
Elements $(k,b)$ first reflect when $b=1$ and then rotate
$k$ times.
Composition $g_1\circ g_2$ means apply $g_2$ first and
$g_1$ second.
The action is faithful and transitive but not free.
Each corner is fixed by the identity and one diagonal
reflection, and the action obeys
$r^4=e$, $s^2=e$, and $srs=r^{-1}$.
This is a vertex-label task, not a task using rendered square
images or full square configurations.
Complete training pairs retain fixed points; the evaluation
definitions below specify which elements and pairs enter
each score.

Each domain uses 40 fixed English template strings with one
state insertion.
Weekday examples include
``The current day is \{state\}.'',
``The calendar shows \{state\}.'', and
``The event happens on \{state\}.''
Corner examples include
``The marked corner is \{state\}.'',
``The square position is \{state\}.'', and
``The diagram label is \{state\}.''
The complete strings are distributed in the dataset definitions.
They are semantically similar short contexts, not independent
natural-language corpora.
FULL contains 280 weekday and 160 corner prompts per model
before partitioning.
Changing the state within a template defines a paired
transformation example.

A seeded permutation of the 40 template IDs selects
28 training, six validation, and six test templates.
Split seeds are 101, 303, 909, 2701, and 4243.
All source and transformed state variants of a template stay
in the same partition.
Across splits, prompts recur; this repeated-partition design
does not create independent datasets.
Validation observations are diagnostic; the reported primary
results use test observations.
Held-out template IDs do not imply disjoint tokenized prefixes
at the state-span readout described below.

\subsection{Models, tokenization, and representation extraction}

The checkpoints are the base releases in the official
\texttt{meta-llama} and \texttt{Qwen} repositories,
not the \texttt{-Instruct} variants.
Model-family sources are the Llama~3 report
\citep{grattafiori2024llama3}, the Llama~3.2 model card
\citep{meta2024llama32}, and the Qwen2.5 report
\citep{qwen2025technical}.
The links in \cref{tab:models} identify the official checkpoint
pages; the revision hashes, rather than the current contents
of those pages, identify the artifacts used in these experiments.

\begin{table}[H]
\centering
\caption{Exact base-model revisions and primary transformer
block output indices.
The embedding hidden state is index zero and is excluded
from FULL.}
\label{tab:models}
\small
\begin{tabular}{lrrrrl}
\toprule
Model & 0.20 & 0.40 & 0.60 & 0.80 & Revision \\
\midrule
\href{https://huggingface.co/meta-llama/Llama-3.2-1B}{Llama-3.2-1B}
& 3&6&10&13
& \texttt{4e20de362430cd3b72f300e6b0f18e50e7166e08}\\
\href{https://huggingface.co/meta-llama/Llama-3.2-3B}{Llama-3.2-3B}
& 6&11&17&22
& \texttt{13afe5124825b4f3751f836b40dafda64c1ed062}\\
\href{https://huggingface.co/meta-llama/Llama-3.1-8B}{Llama-3.1-8B}
& 6&13&19&26
& \texttt{d04e592bb4f6aa9cfee91e2e20afa771667e1d4b}\\
\href{https://huggingface.co/Qwen/Qwen2.5-1.5B}{Qwen2.5-1.5B}
& 6&11&17&22
& \texttt{8faed761d45a263340a0528343f099c05c9a4323}\\
\href{https://huggingface.co/Qwen/Qwen2.5-3B}{Qwen2.5-3B}
& 7&14&22&29
& \texttt{3aab1f1954e9cc14eb9509a215f9e5ca08227a9b}\\
\href{https://huggingface.co/Qwen/Qwen2.5-7B}{Qwen2.5-7B}
& 6&11&17&22
& \texttt{d149729398750b98c0af14eb82c78cfe92750796}\\
\bottomrule
\end{tabular}
\end{table}

The tokenizer revision matches the model revision.
Study~1 uses fast tokenizers, special tokens, right padding,
FP16 model loading, automatic device mapping, and extraction
batches of eight.
Models are in evaluation mode, and extraction uses inference
mode without KV caching.
These are the Study~1 extraction settings; Study~2 uses BF16
model loading as described in \cref{sec:design}.

A character-span search requires exactly one occurrence of
the state string.
Offset mappings select all token positions overlapping that
span, and their hidden vectors are averaged at each selected
layer.
A state span may cover one or more tokens.
Under causal attention, this readout depends on the tokenized
prefix through the last selected token; later tokens do not
affect it when that prefix is unchanged.
The visible state span is part of the input, not an unobserved
answer.
Extraction therefore measures representations of already
presented states.
In Study~1, span means are converted to float32 for storage,
and PCA operates on float64 copies of these readouts.

\paragraph{Training-only PCA.}
For each model, group, layer, and split, PCA is fitted only
to training-template activations.
It subtracts one pooled mean $\hat\mu$ across all training
state--template pairs and uses an SVD basis without whitening.
With orthonormal component rows
$B=U_{\mathrm{PCA}}^\top$, the projected coordinate and
PCA back-projection are
\[
 z=(h-\hat\mu)B^\top,
 \qquad
 \widehat h=zB+\hat\mu.
\]
The back-projection does not recover directions discarded
by truncation.
The same training mean and basis are applied to validation
and test observations.

This pooled training centering differs from subtracting the
state mean separately within each template.
PCA selects directions by total training variation, which
includes both variation among template means and variation
among states within a template.
The total variance retained does not determine how much
within-template state variation is retained.
The diagnostic in \cref{app:template-centering-sensitivity}
operates within the saved 16-dimensional Study~2 subspace
and does not refit PCA.
It compares retaining the original test-bank whitening with
whitening the centered bank by its own second moment.
It does not recover the state variation discarded by the
original projection or replace a PCA fit to full-dimensional
within-template differences.

\paragraph{Dimension settings and sensitivity analyses.}
Study~1 fixed $q=16$ before the final structural run.
The record does not establish that this dimension is optimal
or was selected by a variance-retention threshold.
The recorded sensitivity subset uses dimensions 8, 32,
and 64 at depths 0.4/0.8 and split seeds 101/303.
These sensitivity cells are not pooled with primary results.
A matched summary must restrict the $q=16$ comparator to
the same model, depth, split and endpoint subset;
the archived aggregate table does not impose that restriction
on its $q=16$ row.
\Cref{app:dimension-rank-audit} reports the verified matched
subset.
Additional terminal-layer structural fits in the historical
run are likewise excluded from the primary four-depth means.

These historical prediction and group-law checks do not
establish dimension stability of $\gamma_r$, the exact-class
premium or registration rankings.
Study~2 fitted PCA on its own training templates and retained
16 coordinates for comparability with its development analyses.
Its original spectral analysis used $r=16$.
The subsequent same-bank comparison uses the first eight
saved coordinates and evaluates $r=8$ under its own
whitening metric
(\cref{app:spectral-dimension-sensitivity});
the completed independent-orthogonal comparison is reported
in \cref{app:orthogonal-results}.

\subsection{Independent and joint fitting}

For group element $g$, stack training source rows in $X$
and their matched same-template transformed rows in $Y_g$.
All training states, including fixed-point pairs, are retained.
The no-intercept independent fit solves the unnormalized
ridge objective
\[
 \min_{A\in\R^{q\times q}}
 \left\{
 \|XA-Y_g\|_F^2+0.01\|A\|_F^2
 \right\},
\]
giving
\[
 R_g=(X^\top X+0.01\Id_q)^{-1}X^\top Y_g.
\]
Thus the coefficient $0.01$ refers to a sum-of-squares
objective, not a sample-mean or entry-mean data loss.
The same source design and coefficient are used for every
fitted element.
The operator knows the training source--target correspondence
for that element.
It receives no multiplication table, generator identities,
or relation loss; this is absence of algebraic \emph{loss}
supervision, not absence of transformation supervision.

FULL fits every group element, including the identity.
In particular,
\[
 R_e=(X^\top X+0.01\Id_q)^{-1}X^\top X
\]
is fitted rather than fixed to $\Id_q$.
The common-design ridge identities in \cref{app:ridge-identity}
apply to this fitting contract.

For a matrix $A$, let $\operatorname{mse}(A)$ be its mean
squared entry.
Joint fitting optimizes
\begin{align*}
 L
 &=
 L_{\rm data}+L_{\rm hom}+L_{\rm rel}+L_{\rm id}
 +0.05L_{\rm orth},
 \\
 L_{\rm data}
 &=
 \mathbb E_{(i,t),g}
 \frac1q
 \norm{z_{i,t}R_g-z_{g(i),t}}_2^2,
 \\
 L_{\rm hom}
 &=
 \mathbb E_{g_1,g_2}
 \operatorname{mse}
 (R_{g_2}R_{g_1}-R_{g_1\circ g_2}),
 \\
 L_{\rm id}
 &=
 \operatorname{mse}(R_e-\Id_q),
 \\
 L_{\rm orth}
 &=
 \mathbb E_g
 \operatorname{mse}(R_g^\top R_g-\Id_q).
\end{align*}
The data average includes every training state--template
pair and every group element; the homomorphism average
includes every ordered pair of elements.
For $C_7$,
\[
 L_{\rm rel}
 =
 \operatorname{mse}(R_r^7-\Id_q).
\]
For $D_4$, it is the mean of
\[
 \operatorname{mse}(R_r^4-\Id_q),\qquad
 \operatorname{mse}(R_s^2-\Id_q),\qquad
 \operatorname{mse}(R_sR_rR_s-R_{r^{-1}}).
\]
The fitting losses are not the normalized evaluation metrics.
All elements, including the identity, are jointly optimized
and initialized from ridge.
The full-batch implementation uses float32 parameters and
AdamW with learning rate 0.002, weight decay $10^{-5}$,
1,200 steps, and seed 42.
Exported matrices are converted to float64.
These finite optimization steps do not establish attainment
of a global minimum.
The algebraic terms are soft penalties in the unwhitened
PCA coordinates; they do not enforce membership in the exact
orthogonal-action class defined in the bank's own whitened
metric in \cref{app:spectral-risk}.
One fit per scientific cell is used;
technical optimization seeds are not independent scientific
replications.

Identity predicts $z$ itself.
Random-orthogonal controls use Gaussian matrices followed
by QR with diagonal-sign correction.
For the FULL random-orthogonal control, nonidentity element
matrices are drawn independently and the identity matrix
is fixed exactly.
Its five seeds are 104729, 130363, 155921, 196613, and 262147.
AFS-A instead uses one random-generator family per cell,
generated from a stable cell-specific seed and composed
according to the same words as the focal method.
Neither random baseline is a direct-error-matched null.

The all-identity family satisfies the group laws exactly
and has high prediction error for the nonidentity
transformations in these banks under the Study~1 metric.
Its algebra columns are left unscored by the protocol;
those defects are not mathematically undefined, and zero
algebra error alone would not establish target recovery.
Its Study~1 prediction score is distinct from the
self-whitened identity-family risk $L(I)$
in \cref{sp:identity-reference}.

\subsection{Exact evaluation metrics}
\label{app:metrics}

Let $Z$ contain all held-out source-state rows for one
template, and let $Z_g$ contain the actual same-template
transformed rows in the corresponding source-state order.
For paired matrices $A,B$ with $N$ rows,
\cref{eq:D} defines
\[
 D(A,B)
 =
 \frac1N\sum_{j=1}^N
 \frac{\|A_j-B_j\|_2^2}{\|B_j\|_2^2+10^{-12}}.
\]
Every average below is computed before the documented
higher-level experimental aggregation.
For operator diagnostics, define
\[
 \kappa_{\mathrm{op}}(A,B)
 =
 \frac{\|A-B\|_F}{\|B\|_F+10^{-12}}.
\]
The notation $\kappa_{\mathrm{op}}$ distinguishes this scalar
discrepancy from the orbit-Gram matrix $K$
in \cref{app:spectral-risk}.
The prediction and homomorphism diagnostics are
\begin{align*}
 \Eeq
 &=
 \mathbb E_{g\ne e}D(ZR_g,Z_g),
 \\
 \Ehom
 &=
 \mathbb E_{g_1,g_2}
 D(ZR_{g_2}R_{g_1},ZR_{g_1\circ g_2}),
 \\
 \Ehomop
 &=
 \mathbb E_{g_1,g_2}
 \kappa_{\mathrm{op}}
 (R_{g_2}R_{g_1},R_{g_1\circ g_2}).
\end{align*}
The prediction average excludes the identity element but
retains fixed-point pairs of nonidentity elements.
The homomorphism averages include all ordered element pairs.
Their comparison target is the directly fitted output
$ZR_{g_1\circ g_2}$ or matrix $R_{g_1\circ g_2}$,
not the actual transformed observation $Z_{g_1\circ g_2}$.
In particular, the data-space denominator depends on the
directly fitted prediction.

The metrics $\Ecomp$ and $\Ecompop$ use the same
homomorphism expressions with $g_1\ne e$ and $g_2\ne e$.
Their product may still be the identity.
Their difference from homomorphism is therefore the set of
action pairs, not a separate held-out-element fitting protocol.
Operator diagnostics repeat the same matrix value across
template rows because the fitted matrices are fixed within
a cell.

\paragraph{C7 relation definitions.}
Write $A=R_r$, with $R_k$ the independently available matrix
for $r^k$.
The historical data relation metric is
\[
 \begin{aligned}
 \Erel^{C_7}
 =
 \frac13\Bigg\{&
 D(ZA^7,Z)
 +
 \frac17\sum_{k=0}^6D(ZA^k,ZR_k)
 \\
 &+
 \frac16\sum_{k=1}^6
 D(ZR_kR_{-k\bmod7},Z)
 \Bigg\}.
 \end{aligned}
\]
The historical operator relation metric is only
\[
 \Erelop^{C_7}
 =
 \kappa_{\mathrm{op}}(A^7,R_e).
\]
It is not an unsquared version of the data aggregate.
The $k=0$ term in the data power comparison is
$D(Z,ZR_e)$, whereas the data cycle term compares $ZA^7$
with $Z$.
We retain these distinct comparison targets and frozen metric
definitions rather than redefining a recorded result.

\paragraph{D4 relation definitions.}
Writing $A=R_r$ and $J=R_{r^{-1}}$,
\begin{align*}
 \Erel^{D_4}
 &=
 \tfrac13\{
 D(ZA^4,Z)
 +D(ZR_s^2,Z)
 +D(ZR_sAR_s,ZJ)
 \},
 \\
 \Erelop^{D_4}
 &=
 \tfrac13\{
 \kappa_{\mathrm{op}}(A^4,R_e)
 +\kappa_{\mathrm{op}}(R_s^2,R_e)
 \\
 &\hspace{3.5em}
 +\kappa_{\mathrm{op}}(R_sAR_s,J)
 \}.
\end{align*}
The data definition compares the first two relations with $Z$,
while the operator definition compares with the learned $R_e$.
Relation components are also retained individually.
All data metrics in these relation definitions are squared
relative per-row errors; all $\kappa_{\mathrm{op}}$
diagnostics are unsquared.
A lower raw data number than operator number does not by itself
show a quantitative gap between supported and unsupported
directions.

\paragraph{Auxiliary inverse metrics.}
For a fixed template, the historical inverse data metric is
\[
 \frac1{|G|}\sum_{g\in G}
 \frac{
 \mathbb E_i
 \|z_iR_gR_{g^{-1}}-z_i\|_2^2
 }{
 \mathbb E_i\|z_i\|_2^2+10^{-12}
 }.
\]
This is a ratio of means within each element, unlike $D$'s
mean of ratios.
Its operator counterpart is
\[
 \frac1{|G|}\sum_{g\in G}
 \frac{\|R_gR_{g^{-1}}-\Id_q\|_F^2}{q}.
\]
Both inverse metrics are squared and average all group
elements, including the identity.
They therefore differ from both the unsquared
$\kappa_{\mathrm{op}}$ diagnostics and the nonidentity
inverse component inside $\Erel^{C_7}$.

\paragraph{FULL word diagnostics.}
For a word $w=(a_1,\ldots,a_k)$ listed in application order,
the matrix product is $R_{a_1}\cdots R_{a_k}$ and the target
element is $a_k\circ\cdots\circ a_1$.
FULL compares these products with actual transformed target
rows.
For $C_7$, the word set is $r^k$, $k=1,\ldots,6$,
giving six words.
For $D_4$, it contains all ordered words over $\{r,s\}$
of lengths one through three, giving 14 words.
Each word receives equal weight.
Words with repeated group-element endpoints are retained,
including D4 words that return to the identity; their
target is the actual source observation, not $ZR_e$.

The parent FULL operator systems were trained on all elements,
and Joint can also use all relations.
These word scores therefore include already fitted elements
and are not generator-only training tests or uniform averages
over distinct group elements.
AFS-A separately withholds nongenerator element supervision
and uses the evaluation set specified below.

\subsection{AFS-A/B/C controls}

\paragraph{AFS-A: generator-only fitting.}
The focal AFS-A fit uses training-only PCA and fits only $r$
for $C_7$ or $r,s$ for $D_4$, with the same ridge coefficient.
C7 evaluates $r^2,\ldots,r^6$;
$r^7$ is a separate cycle diagnostic.
D4 uses breadth-first shortest generator words for the five
elements in $G\setminus\{e,r,s\}$.
Matrix products follow left-to-right row-vector application
order, and the target element is obtained by the corresponding
action lookup.

Controls are identity, random generator products,
directly fitted element-specific ridge, and D4 matched-length
words that evaluate to a different element.
All focal errors compare with actual transformed test
activations.
``Unseen'' refers to the excluded operator elements,
not new state labels.
Evaluation also holds out templates.
Different D4 elements can act identically on an individual
corner; the wrong-word contrast averages over the complete
state set.
Word length is matched, while individual source-to-target
distance is not an additional matched condition.

\paragraph{AFS-B: group-element label permutations.}
AFS-B retains the historical independent ridge operators and
PCA coordinates.
It enumerates identity-preserving permutations of
group-element labels, excludes group automorphisms,
and draws 5,000 non-automorphic permutations with replacement
per cell.
The same multiplication table is evaluated after relabeling.
We retain the recorded smoothed lower-tail score,
\[
 s_{\rm MC}
 =
 \frac{1+\#\{b:E_b\le E_{\rm true}\}}{5001}.
\]
The observed true registration is not established as
exchangeable with these restricted draws, which exclude
automorphisms including the identity registration.
Accordingly, this score describes separation from the chosen
reference ensemble; no calibrated $p$-value, Type-I-error
control or semantic-null rejection is claimed.
Historical files retain their original
\texttt{empirical\_p} field name and values.
Repeated cells share representations, model families,
and templates.
A universal pass count is not a count of independent discoveries.

AFS-B fixes the fitted operators and permutes their
group-element labels.
These non-automorphic relabelings are distinct from the
coherent alternative state-action registrations and
permutation-image comparisons used in Study~2
(\cref{sec:registration-sensitivity}).
They do not define the same comparator class or evaluate
the same spectral criterion.

\paragraph{AFS-C: corrupted training correspondences.}
AFS-C retains the original activation vectors,
template membership and training-fitted PCA.
Let $z_{i,t}$ denote the coordinate of state $i$ in template $t$.
In the state-relabel condition, each training template receives
a seeded permutation $\pi_t$, and the training pairs become
\[
 z_{\pi_t(i),t}
 \longmapsto
 z_{\pi_t(g(i)),t}.
\]
The same $\pi_t$ is used for all elements within that template,
so the induced state action is the coherent conjugate
$\pi_t\circ g\circ\pi_t^{-1}$.
Permutations are sampled separately by training template;
no single global relabeling is imposed.
Permutations commuting with every registered action are excluded.

In the target-scramble condition, the source remains $z_{i,t}$
and the target becomes $z_{\sigma_{t,g}(g(i)),t}$.
For each training template and nonidentity element,
a nonidentity target permutation $\sigma_{t,g}$ is sampled
separately.
Identity-action pairs remain self-pairs, and individual
fixed points are allowed.
All element-specific Ridge maps are refitted with the
original regularization.
Direct prediction is evaluated against the unchanged true
validation/test transformations.
The algebraic diagnostics retain their historical definitions,
including comparisons with condition-specific fitted outputs;
they are not additional semantic-target accuracy measurements.

The five control seeds are 70001, 70003, 70009, 70019 and 70039.
Neither condition fits one globally coherent alternative
and evaluates it against consistently relabeled held-out targets.
The sharing of these template-level permutations across
models, layers and splits is described in \cref{app:inference}.

\subsection{Rank-matched lexical comparison}
\label{app:lexical-rank}

The embedding readout averages token embeddings overlapping
the unique target span; it does not average contexts.
The saved arrays contain exactly seven distinct weekday
vectors or four distinct corner vectors, with each state's
vector bitwise identical across its 40 templates.
The audit selects one observed row per state only after
verifying this equality.
Centering $n$ unique rows gives the upper bound
$\min(d,n-1)$, not an automatic rank equality.
Float64 SVD diagnostics confirm centered ranks 6 and 3
for the four archived lexical inputs, and reproduce all
20 recorded training-split rank checks
(\cref{app:dimension-rank-audit}).

The original $q=16$ lexical failure remains in the record.
The separate rank-matched protocol uses $q=6$ and $q=3$
for both lexical and internal readouts;
matching refers to the ranks of their training projected
representations.
It covers Llama-3.2-1B and Qwen2.5-1.5B, both groups,
five partitions, and four internal depth slots:
20 embedding and 80 internal fit cells.
Embedding results have no contextual depth axis and are
reused when forming depth comparisons.
The later replay reused the same embedding caches and rank
records; it supplies no independent lexical observations.

The numerical AFS-D table retains the recorded independent-fit
results at these dimensions.
It is not a dimension-16 embedding comparison and cannot be
merged numerically with the primary six-model FULL result.
All 80 original internal training inputs have sufficient rank
for their declared truncation:
their centered ranks are 195 for weekdays and 111 for corners.
The recorded PCA procedure retains the leading $q$ right
singular vectors.
Together with the audited nonzero retained singular values,
this implies training-coordinate rank $q$ at both readouts,
as shown in \cref{app:dimension-rank-audit}.

The historical record retains evaluation metrics but not
the fitted PCA bases, coordinates or operator matrices
for this lexical comparison.
Thus the training-rank match is supported by the inputs and
algorithm, not by direct remeasurement of those historical
outputs.
Training rank alone does not establish the rank of every
held-out internal bank or any fitted operator.

The verified affine rank supports the finite-orbit construction
in \cref{prop:finite} on each lexical input's full centered
$(n-1)$-dimensional span.
This establishes exact linear realizability of every state
permutation on that span.
Orthogonality additionally follows after uniform-state
second-moment whitening of the full span
(\cref{app:finite-orbit-proof}).
The historical PCA coordinates were not whitened,
and the regularized fitted errors do not certify exactness
or orthogonality.
Nor do these ranks show that embedding training discovered
group-specific algebra.

Matching training projected rank does not match retained
information.
The leading $q$ directions explain 54.42--69.83\% of the
original internal training variance for weekdays and
45.81--73.16\% for corners, whereas the lexical input span
has only 6 or 3 dimensions.
The internal fractions use the full original training spectrum,
not the retained coordinates as denominator.
They measure total variance retention, not retention of
within-template state variation.
The later Study~2 centering diagnostic neither reconstructs
the missing historical outputs nor measures the full-dimensional
state variation discarded in these lexical/internal comparisons.
Internal-minus-embedding intervals remain historical matched
comparisons with the fixed-panel and repeated-partition
limitations.

\subsection{Inference, aggregation, and protocol scope}
\label{app:inference}

The supplied summaries preserve the completed experiments'
estimands and rounding.
Primary FULL means select dimension 16, four specified depths,
five split seeds, test templates, and the fixed six-model panel.
The FULL random-orthogonal comparator additionally averages
its five technical control seeds.
AFS-A instead uses the single cell-specific random-generator
family described above.
Paired contrasts match method observations on model, group,
depth/layer, dimension, split, template, and relevant
intervention fields.
Exact duplicate rows differing only in technical optimization
labels are collapsed.
For a FULL random comparator, the single learned estimate
is matched against each random seed, which remains a
resampling level for that contrast.

The FULL hierarchy is
\[
 \begin{gathered}
 \text{model}\longrightarrow\text{depth slot}
 \longrightarrow\text{split seed}
 \\
 \longrightarrow
 \text{random-baseline seed, when applicable}
 \longrightarrow\text{template}.
 \end{gathered}
\]
Each nested level is sampled with replacement at its observed
size; terminal observations are concatenated and averaged.
Percentile endpoints use 5,000 draws.
The base seed is 314159, with the frozen analysis- and
metric-specific offsets recorded in the runner;
the base seed alone does not reproduce every interval.

AFS-A resamples model, depth, split and template.
Each resampled template carries all five excluded-element
measurements together.
These element measurements are retained as terminal
observations, not independently resampled as another
hierarchy level.
AFS-C additionally resamples the scramble seed.
AFS-A and AFS-C use bootstrap seed 314159 directly,
without the offsets used by some FULL analyses.
The AFS-C hierarchy is model, depth, split, scramble seed
and template, with each real observation paired to all five
corresponding scramble fits.
The five excluded elements in AFS-A and the five scramble
seeds in AFS-C are different kinds of repeated observations.

AFS-D instead holds its two models and depth slots fixed
within each reported model/group contrast, resampling split
and template for data metrics and split only for operator
metrics.
The global additive-field analysis first forms paired,
equally weighted model/depth and group aggregates within
each split, then resamples the five split-level contrasts.
It does not resample five independent datasets or use
the FULL hierarchy.

These procedures retain their recorded observation weighting
rather than replacing the FULL mean.
For the direct-error aggregate reported in
\cref{app:results}, the additive-field analysis gives
equal weight to domains and model--depth conditions
within each split.
Its Ridge and Joint point estimates coincide with
the equal-domain averages of the primary FULL means,
and its global additive point estimate coincides
with the corresponding identity baseline.
The uncertainty procedures remain distinct:
the additive-field analysis resamples five paired
split-level contrasts, whereas FULL uses the
hierarchy described above.
The original seed stream and aggregation code accompany
each table.

This hierarchy is the implementation's historical procedure,
not a proof that all its axes are exchangeable independent
sampling units.
Models are selected from two related families, depths are
repeated locations in the same model, and different split
seeds repeatedly partition the same templates.
Resampling lower levels separately does not model every
crossed dependency or remove reuse of templates across
partitions.
Accordingly, the intervals quantify this resampling calculation;
we do not claim calibrated population coverage over unseen
model families or independently sampled datasets.
The split variability is best interpreted as robustness to
the recorded partitions.
Technical seed 42 adds no replication.

AFS-C's permutation keys include the control seed, variant,
domain and template, and additionally the element for target
scrambling; they do not include model, depth or split.
The same template-level assignments are therefore shared
across models, layers and overlapping partitions.
The nested bootstrap does not encode all these crossed
dependencies.
Its intervals retain their role as sensitivity summaries
under the recorded resampling scheme, rather than calibrated
population coverage or a semantic-specificity test.

These historical intervals are not applications of the
finite-context certificate in \cref{sp:certificate}.
The recorded template banks and resampling hierarchies do
not supply its stipulated independent population context
sample and fixed population norm bound.
Nor do these resampling intervals change the deterministic
finite-bank interpretation of the spectral analysis
in \cref{sp:finite-bank}.

Protocol files and source/configuration hashes document frozen
analysis choices and completed run lineage.
A frozen local protocol is not equivalent to an independently
timestamped public preregistration.
Only the gate family identified by the original frozen local
protocol was designated confirmatory within that run.
This designation does not establish independent public
preregistration, calibrated Type-I-error control or global
multiplicity correction.
Other contrasts, post-hoc diagnostics, and the analytical
observations in this paper should not be read as globally
multiplicity-corrected confirmatory tests.
These historical sections preserve the original evidence;
the subsequent coordinate study and its diagnostics are
reported separately and do not retroactively upgrade a
historical gate or supply missing native-intervention outcomes.
